\scanpage{5}
% The scan is upside down; text is transcribed in its own reading orientation.
\mspage{936 \deleted{quater} \handwritten{(iv)}}
% archive sheet number: 4

tions de texte (par rapport à l'édition provisoire d'Octobre dernier), chaque fois que j'ai trouvé que c'était plus simple que d'ajouter des notes de bas de page correctrices, et que cela ne modifiait pas de façon draconienne le sens et l'esprit du texte initial. J'ai enfin supprimé un passage d'une page et demie dans la sous-note ``Premiers ennuis - ou les caïds d'outre-Pacifique'' (\deleted{\unsure{In}} in ``La maffia'' n° $171_2$, (b)),\deleted{\unsure{mais pour}} en les remplaçant par un court commentaire rétrospectif.

\noindent\rule{2cm}{0.4pt}

Il est heureux que je n'aie pas suivi Mebkhout\deleted{\unsure{,}} pour donner à son contentieux avec Kashiwara une place qui ne lui revenait pas, dans \deleted{\unsure{le}} \handwritten{\inserted{mon}} compte-rendu de la mystification collective \inserted{sans précédent} qui a eu lieu autour de son oeuvre. Aussi je pense que les pages qui suivent, malgré les passages incorrects ou tendancieux (signalés en leur lieu) concernant le rôle et les mérites de Mebkhout, \deleted{\unsure{donnent|demeurent}} sont \handwritten{\inserted{pourtant}} un reflet circonstancié et fidèle d'un certain ``esprit du temps'' tel qu'il règne à présent (entre autres, sûrement) dans \deleted{\unsure{les}} \inserted{certains} milieux mathématiques, y compris parmi les plus huppés. \handwritten{\unsure{(}}Et ces deformations même dont je m'étais fait l'écho et la caution, m'apparaissent elles aussi comme un signe éloquent de cet esprit, dont mon ami Zoghman Mebkhout\deleted{\unsure{,}} \deleted{arrivé sur le tard} (j'ai fini par m'en rendre compte) \deleted{enfin)} est lui aussi,\deleted{un des produits} dans certains de ses réflexes et comportements, un produit.

\handwritten{(}171 (i)\handwritten{)}\quad L'ancêtre\quad (*)

(1 mars et 2-8 mai (1985)) \deleted{\unsure{(i)}} \deleted{\unsure{D}}\handwritten{D}ans chacune de ces ``quatre opérations'' partielles que \handwritten{\inserted{\unsure{j'ai}}} distinguées dans mon enterrement anticipé, c'est Deligne qui visiblement joue le rôle de chef d'orchestre (ou plutôt, de Grand Officiant aux

\noindent\rule{2cm}{0.4pt}

(*\deleted{\unsure{*}}) (13 mai) La présente note et les quatre suivantes formaient initialement une note unique, ``L'Apothéose'' (n° 171), du 1 mars. Elle incluait également la note précédente ``Les joyaux'' (n° 170(iii)). Elle a été reprise et considérablement étoffée entre les 2 et 8 mai, surtout en ce qui concerne la partie mathématique, et scindée en les quatre notes séparées ``L'ancêtre'', ``L'oeuvre...'', ``... et l'aubaine'', ``Le jour de gloire'' (n°s 171 (i) à (iv)), en plus de la note ``Les joyaux'' déjà mentionnée. Il s'y ajoutent les huit sous-notes (n°171 (v) à (xii)) se rapportant aux quatre notes en question, et les quatre sous-notes (n° $171_1$ à $171_4$) du mois d'avril, faisant le récit des étranges mésaventures de
